
To conclude this section, we examine the case in which \(G\) is
``semisimple'' over \(S\).

\paragraph{\textbf{Theorem 3.18.}}
\SGARefTarget{sga3:xiv:3.18}
{\itshape Let \(S\) be a prescheme and let \(G\) be a smooth
\(S\)-prescheme in groups whose geometric fibers are semisimple
``adjoint'' groups, that is, semisimple groups whose reductive
center \textup{(\SGARefLink{sga3:xii:4.1}{XII~4.1} and \SGARefLink{sga3:xiv:4.4}{4.4})} is the identity group. Then the
subgroups of type \((\mathrm C)\) of \(G\) are precisely its
maximal tori, and hence also precisely its Cartan subgroups
\textup{(\SGARefLink{sga3:xii:3.1}{XII~3.1})}.}

In view of the definitions, it is enough to treat the case where
\(S\) is the spectrum of an algebraically closed field. We must
then show that, for a maximal torus \(T\) of \(G\), its Lie algebra
\(\mathfrak t\) is a Cartan subalgebra of \(\mathfrak g\). Put
\(r=\dim T=\operatorname{rank}\mathfrak t\). Taking into account
the inequality
\[
  \operatorname{nil\,rank}(\mathfrak g)
  \geq
  \operatorname{nil\,rank}(G)
  =
  \dim T
  =
  \operatorname{rank}\mathfrak t
  =
  r,
\]
it is enough to find \(x\in\mathfrak t\) such that
\[
  \dim\Nil(x,\mathfrak g)=r.
\]
Since \(\mathfrak t\) is abelian, and therefore nilpotent, this is
equivalent to finding \(a\in\mathfrak t\) for which
\[
  \operatorname{ad}(a)_{\mathfrak g/\mathfrak t}
\]
is injective \textup{(\SGARefLink{sga3:xiii:corollary:5:7}{XIII~5.7}\textup{(a)})}.

Consider the characters \(\alpha\) of \(T\) occurring in the
representation of \(T\) induced by the adjoint representation of
\(G\). The structure theory of the semisimple group \(G\)
\textup{\emph{(Bible)}, 13, \SGARefLink{sga3:xiv:7:lemma1}{Theorem~1}\textup{(a)}, and Theorem~3,
\SGARefLink{sga3:xiv:7:lemma2}{Corollary~2}}, or more precisely the ``big cell'' of \(G\), shows
that the eigenspace of \(\mathfrak g\) for the trivial character is
exactly \(\mathfrak t\). The other eigenspaces are
one-dimensional, and their characters are precisely the roots of
\(G\) relative to \(T\). Here the big cell is built from \(T\) and
the subgroups \(P_\alpha\), each isomorphic to the additive group
\(\mathbf G_a\), stable under \(T\), and corresponding to the root
characters.

The reductive center of \(G\) is the intersection of the kernels
of the characters of \(T\) occurring in the adjoint representation
\textup{(\SGARefLink{sga3:xii:4.8}{XII~4.8})}. Thus the assertion that \(G\) is adjoint says
that its roots generate the lattice
\[
  M=\operatorname{Hom}(T,\mathbf G_m).
\]
A familiar lemma from root theory says that every root belongs to
a system of simple roots, hence to a basis of the group generated
by the roots, and therefore to a basis of the dual lattice of
\(T\). We obtain the following.

\paragraph{\textbf{Corollary 3.19.}}
\SGARefTarget{sga3:xiv:3.19}
{\itshape Let \(G\) be an adjoint semisimple algebraic group over
an algebraically closed field \(k\), and let \(T\) be a maximal
torus of \(G\). For every root
\[
  \alpha\colon T\longrightarrow\mathbf G_m
\]
of \(G\) relative to \(T\), the corresponding homomorphism
\[
  \alpha'\colon\mathfrak t\longrightarrow k
\]
is nonzero.}

This assertion is essentially equivalent to \SGARefLink{sga3:xiv:3.18}{Theorem 3.18}. Indeed,
for \(t\in\mathfrak t\), the endomorphism
\(\operatorname{ad}(t)\) is semisimple, and its eigenvalues on
\(\mathfrak g/\mathfrak t\) are precisely the numbers
\(\alpha'(t)\). Hence
\(\operatorname{ad}(t)_{\mathfrak g/\mathfrak t}\) is injective if
and only if all the \(\alpha'(t)\) are nonzero; and such a
\(t\in\mathfrak t\) exists if and only if every \(\alpha'\) is
nonzero.
